\documentclass[a4paper]{article}
\usepackage{xcolor,geometry,graphicx,palatino}
\addtolength{\parskip}{0.5\baselineskip}
\setlength{\parindent}{0pt}

\title{Notes on Linear Algebra}
\author{Rich Wareham}

\newcommand{\hmatrixrule}{\rule[0.5ex]{2em}{0.7pt}}
\newcommand{\vmatrixrule}{\rule[-0.5ex]{0.7pt}{1.5em}}
\newcommand{\highlight}[1]{{\color{blue}#1}}

\begin{document}
\maketitle

\section{Matrix vector multiplication}

I find it confusing to think about vectors and matrices as seas of numbers.
Instead I find it simpler, and more geometrically intuitive, to see them as
stacked sets of row or column vectors. This is the approach I'll use in these
notes and is also one I suggest you attempt to master. It can often reduce a
linear algebra proof to a few lines.

\subsection{Left-multiplication}

A $1 \times N$ vector, $\vec{x}$, multiplied by a $N \times M$ matrix, $A$,
gives a $1 \times M$ vector which is the dot product of $\vec{x}$ with the
\emph{columns} of $A$.

\[
\begin{array}{cccc}
\left[
  \begin{array}{cccc}
    \hmatrixrule & \vec{x} & \hmatrixrule 
  \end{array} \right] &
\left[
	\begin{array}{cccc}
		\vmatrixrule & \vmatrixrule &  & \vmatrixrule \\[1ex]
		\vec{a}_1 & \vec{a}_2 & \cdots & \vec{a}_M \\[1ex]
		\vmatrixrule & \vmatrixrule &  & \vmatrixrule 
	\end{array}
\right] &
= &
\left[ \begin{array}{cccc}
	\vec{x} \cdot \vec{a}_1 & \vec{x} \cdot \vec{a}_2 & \cdots & \vec{x} \cdot \vec{a}_M
      \end{array} \right] 
\\
\rule{0pt}{2em} 1 \times N & N \times M & \rightarrow & 1 \times M
\end{array}
\]

\subsection{Right-multiplication}

A $N \times M$ matrix, $A$, multiplied by a $M \times 1$ vector, $\vec{x}$,
gives a $N \times 1$ vector which is the dot product of $\vec{x}$ with the
\emph{rows} of $A$.

\[
\begin{array}{cccc}
\left[
	\begin{array}{ccc}
		\hmatrixrule & \vec{a}_1 & \hmatrixrule \\
		\hmatrixrule & \vec{a}_2 & \hmatrixrule \\
		            & \vdots & \\
		\hmatrixrule & \vec{a}_N & \hmatrixrule
	\end{array}
\right] &
\left[
	\begin{array}{c}
		\vmatrixrule \\[1ex]
		\vec{x} \\[1ex]
		\vmatrixrule 
	\end{array}
\right] &
= &
\left[
	\begin{array}{c}
		\vec{a}_1 \cdot \vec{x} \\ \vec{a}_2 \cdot \vec{x} \\ \vdots \\ \vec{a}_N \cdot \vec{x}
	\end{array}
\right]
\\
\rule{0pt}{2em} N \times M & M \times 1 & \rightarrow & N \times 1
\end{array}
\]

\subsection{The four fundamental subspaces}

\subsubsection{The row space}

The row space of a $N \times M$ matrix, $A$, is \highlight{the $M$-component space
spanned by the \emph{rows} of $A$}. If $\vec{x}$ is in the row space of $A$
then at least one of $\vec{x} \cdot \vec{a}_i$, where $\vec{a}_i$ is a row of
$A$, \emph{must} be non-zero by definition. \highlight{The \emph{dimension} of the row
space is equal to the rank of the matrix.}

\subsubsection{The null space}

The null space of a $N \times M$ matrix, $A$, is the \highlight{$M$-component} space of
vectors, $\vec{x}$, for which \highlight{$A\vec{x} = \vec{0}$}. Looking at our
definition of the row space, it is clear that if a vector is in the row space
it \emph{cannot} be in the null space and \emph{vice versa}. Therefore the row
space is the \emph{compliment} of the row space.

\subsubsection{The column space}

The column space of a $N \times M$ matrix, $A$, is \highlight{the $N$-component
space spanned by the \emph{columns} of $A$}. If $\vec{x}$ is in the column
space of $A$ then at least one of $\vec{x} \cdot \vec{a}_i$, where $\vec{a}_i$
is a column of $A$, \emph{must} be non-zero by definition. \highlight{The
\emph{dimension} of the column space is equal to the rank of the matrix.}

\subsubsection{The left null space}

The left null space of a $N \times M$ matrix, $A$, is the \highlight{$N$-component}
space of vectors, $\vec{x}$, for which \highlight{$\vec{x}^TA = \vec{0}^T$}. Looking
at our definition of the column space, it is clear that if a vector is in the
column space it \emph{cannot} be in the left null space and \emph{vice versa}.
Therefore the column space is the \emph{compliment} of the left null space.

\subsubsection{The rank of a matrix}

The rank of a matrix is an important concept. As we have seen, if a matrix is
not \emph{full-rank} then the left-null space and null space must be non-zero
and so some information from a vector will always be `thrown away' (or set to
zero) by the matrix.

\section{Matrix-matrix multiplication}

Like matrix-vector multiplication, it is best to think of this as a set of
row-wise and column-wise operations.

\[
\begin{array}{cccc}
\left[
	\begin{array}{ccc}
		\hmatrixrule & \vec{a}_1 & \hmatrixrule \\
		\hmatrixrule & \vec{a}_2 & \hmatrixrule \\
		            & \vdots & \\
		\hmatrixrule & \vec{a}_N & \hmatrixrule
	\end{array}
\right] &
\left[
	\begin{array}{cccc}
		\vmatrixrule & \vmatrixrule &  & \vmatrixrule \\[1ex]
		\vec{b}_1 & \vec{b}_2 & \cdots & \vec{b}_M \\[1ex]
		\vmatrixrule & \vmatrixrule &  & \vmatrixrule 
	\end{array}
\right] &
= &
\left[
	\begin{array}{cccc}
\vec{a}_1 \cdot \vec{b}_1 & \vec{a}_1 \cdot \vec{b}_2 & \cdots & \vec{a}_1 \cdot \vec{b}_M \\
\vec{a}_2 \cdot \vec{b}_1 & \vec{a}_2 \cdot \vec{b}_2 & \cdots & \vec{a}_2 \cdot \vec{b}_M \\
\vdots & \vdots & \ddots & \vdots \\
\vec{a}_N \cdot \vec{b}_1 & \vec{a}_N \cdot \vec{b}_2 & \cdots & \vec{a}_N \cdot \vec{b}_M 
	\end{array}
\right]
\\
\rule{0pt}{2em} N \times L & L \times M & \rightarrow & N \times M
\end{array}
\]

\subsection{Orthonormal matrices}

An orthonormal matrix is one where all columns are \highlight{independent} and
\highlight{have unit length}. If $\{ \vec{a}_i : i \in 1 \dots M \}$ are the
columns of an orthonormal matrix then it follows from the definition that
\[\highlight{
  \vec{a}_i \cdot \vec{a}_i = 1, \quad \vec{a}_i \cdot \vec{a}_j = 0 \qquad
  \forall \  i, j \in 1 \dots M.
}\]
By using this result and looking at our representation of matrix-matrix
multiplication above it should be clear that, if $Q$ is an orthonormal matrix,
$Q^TQ = I$ and hence \highlight{the transpose of an orthonormal matrix is its
inverse}.

\subsubsection{The rows of $Q$}

Further, if $Q^T$ is the inverse of $Q$, then $QQ^T = I$ and hence $Q^T$ must
itself be orthonormal. Since, by definition of an orthonormal matrix, the columns
of $Q^T$ are independent and have unit length then it follows that
\highlight{the rows of an orthonormal matrix are \emph{also} independent and
unit length}.

\subsubsection{The geometric interpretation of orthonormal matrices}

The geometric interpretation of the action of an orthonormal matrix can easily
be seen by considering the actions of an $N \times M$ orthonormal matrix $Q$
with \emph{rows} $\{ \vec{q}_i : i \in 1 \dots N \}$ on a $M \times 1$ vector
$\vec{x}$. Each component of $Q\vec{x}$ is $\vec{x}$ \emph{resolved} onto one of
the rows of $Q$. Since the rows of $Q$ are independent and have unit length,
this results in a change of co-ordinates for $\vec{x}$, resolving it onto the
basis formed by the rows of $Q$.  In summary, \highlight{an $N \times M$
orthonormal matrix resolved a $M$-component vector into a $N$-dimensional
subspace} and \highlight{a square orthonormal matrix transforms from one
co-ordinate system to another without a loss of information}. This last
statement also justifies our implicit assertion above that all square
orthonormal matrices are invertible.

\section{Eigenvalues and eigenvectors}

An \emph{eigenvector}, $\vec{v}$, of some \highlight{square matrix}, $A$, is
defined to be any vector for which \highlight{$A \vec{v} = \lambda \vec{v}$}. By
\emph{convention} we choose that \highlight{eigenvectors have unit length}
although we are in general free to choose the length of eigenvectors. The value
$\lambda$ is an \emph{eigenvalue} of the matrix, $A$.  \highlight{We do not
consider the vector $\vec{0}$ to be an eigenvector} of a matrix since it
trivially satisfies the requirements for all square matrices.

A \emph{left-eigenvector} of some square matrix, $A$, is a vector, $\vec{v}$,
which satisfies \highlight{$\vec{v}^T A = \lambda \vec{v}^T$}. It is obvious
that \highlight{the eigenvectors of a matrix are the left-eigenvectors of its
transpose}.

\subsection{Eigenvectors of symmetric matrices}

In general eigenvectors need not be orthogonal to each other but there is a
special case where they are. Suppose the square matrix, $A$, is symmetric so
that $A^T = A$. In this case the left-eigenvectors and left-eigenvalues are the
same as the eigenvectors and eigenvalues.

Suppose two eigenvectors, $\vec{e}_1$ and $\vec{e}_2$ with associated
eigenvalues $\lambda_1$ and $\lambda_2$, were non-orthogonal. In which case, we
could represent one as some offset from the other: $\vec{e}_2 = \alpha \vec{e}_1
+ \vec{\Delta}$ where $\vec{e}_1 \cdot \vec{\Delta} = 0$.  Now consider
multiplying $A$ on the right by $\vec{e}_2$:

\begin{centering}\begin{tabular}{rl}
\parbox{0.4\columnwidth}{\centering
\begin{eqnarray}
  A \vec{e}_2 &=& \lambda_2 \vec{e}_2 \nonumber \\
  &=& \lambda_2 \alpha \vec{e}_1 + \lambda_2 \vec{\Delta} \label{eqn:res-1}
\end{eqnarray}\\
}&
\parbox{0.4\columnwidth}{\centering
\begin{eqnarray}
  A \vec{e}_2 &=& A \vec{e}_1 + A \vec{\Delta} \nonumber \\
  &=& \lambda_1 \vec{e}_1 + A \vec{\Delta} \label{eqn:res-2}
\end{eqnarray}\\
}\\
\end{tabular}
\end{centering}

\noindent By (\ref{eqn:res-1}), (\ref{eqn:res-2}) and the orthogonality of
$\vec{e}_1$ and $\vec{\Delta}$ it follows that $A\vec{\Delta} = \lambda_2
\vec{\Delta}$.  Hence, by definition, $\vec{e}_2 = \vec{\Delta}$ and is
orthogonal to $\vec{e}_1$.

This is a sketch proof and is non-rigorous but provides a justification for the
claim that \highlight{the eigenvectors of a symmetric matrix are orthogonal}. A
rigorous proof adds the condition that the matrix be of a form known as positive
semi-definite but this is beyond the scope of the course.

\subsection{Relation to the fundamental subspaces of a matrix}

Looking at the definition of right-multiplication above it should be clear that
if a \emph{non-zero} vector $\vec{b}$ can be expressed via another vector
$\vec{x}$ applied to $A$ so that $A\vec{x} = \vec{b}$ then $\vec{x}$ is in the
row space of the matrix. In other words, $\vec{x}$ is non-orthogonal to at
least one row of $A$ since $\vec{b}$ has at least one non-zero component.a
By a similar argument, if $\vec{x}^T A = \vec{b}^T$ and $\vec{b}$ is non-zero,
then $\vec{x}$ must be in the \emph{column space} of $A$.

It follows that \highlight{all eigenvectors of a square matrix must be in the
row space} and \highlight{all left-eigenvectors of a square matrix must be in
the column space}. Since, for a symmetric matrix, the left-eigenvectors are the
same as the eigenvectors, \highlight{for a symmetric matrix, the eigenvectors
must lie in \emph{both} the column space and the row space}.

As the eigenvectors for a symmetric matrix must be orthogonal and are by
convention unit length, it may be no surprise to you that \highlight{the
eigenvectors of a symmetric matrix form an orthonormal basis for the row and
column spaces}. It is instructive to attempt to show this. It may be done quite
simply using a similar method to that used to show the orthogonality of
a symmetric matrix's eigenvectors above: consider some vector composed of
multiples of the eigenvectors and a remainder term orthogonal to all
eigenvectors and then right-multiply the matrix by it.

Since the eigenvectors of a symmetric basis span the column and row spaces and
since they are orthonormal, it should be cleat that \highlight{the number of
eigenvectors for a symmetric matrix equals its rank}.

\subsection{The eigenvector decomposition}

Consider some $N \times M$ matrix, $A$. We can form two symmetric matrices from
it: the $N \times N$ matrix $AA^T$ and the $M \times M$ matrix $A^TA$. If $V$ is
a matrix whose columns are eigenvectors of $A^TA$ and $U$ is a matrix whose
columns are eigenvectors of $AA^T$ then, by definition,
\[
  A^TA V = V \Lambda_V, \qquad AA^T U = U \Lambda_U,
\]
where $\Lambda_V$ and $\Lambda_U$ are matrices whose diagonal elements
containing the appropriate eigenvalues. The size of $\Lambda_U$ and $\Lambda_V$
depend on the \emph{ranks} of $AA^T$ and $A^TA$ respectively. If you look up the
properties of ranks, you'll discover that the ranks of $A^TA$ and $AA^T$ are
equal to the rank of $A$. Let's call it $R$. In this case, therefore,
\highlight{the matrices $\Lambda_U$ and $\Lambda_V$ are $R \times R$, where $R =
\mbox{rank}(A)$}.

Since we know that $V$ and $U$ are orthonormal, their inverses must be their own
transpose and so we can rearrange the above as follows:
\[
  A^TA = V \Lambda_V V^T, \qquad AA^T = U \Lambda_U U^T.
\]
This is called the \emph{eigenvector decomposition} of $A$. 

\subsubsection{Eigenvalues of $AA^T$ and $A^TA$}

Imagine, for the moment, that $\vec{v}$ is an eigenvector of $A^TA$ with
associated eigenvalue $\lambda_v$.
Then, by the definition,
\[
  A^TA\vec{v} = \lambda_v \vec{v}
  \quad \Rightarrow \quad
  AA^TA\vec{v} = \lambda_v A\vec{v}.
\]
Hence if $\vec{v}$ is an eigenvector or $A^TA$, then $A\vec{v}$ is an
eigenvector of $AA^T$ with the same eigenvalue. Similarly, if $\vec{u}$ is an
eigenvector of $AA^T$ then $A^T\vec{u}$ is an eigenvector of $A^TA$ with the
same eigenvalue. In summary, \highlight{the eigenvalues of $A^TA$ and $AA^T$ are
identical}.

Because of this, we can always \emph{choose} the ordering of $U$ and $V$ so as
to make the diagonal terms of $\Lambda_U$ and $\Lambda_V$ identical and hence
make both matrices equal to some diagonal matrix $\Lambda$. The way this is done
is that, conventionally, \highlight{the columns of $U$ and $V$ are ordered by
decreasing eigenvalue}.

In summary, when the columns of $U$ and $V$ are arranged in decreasing order of
eigenvalue, we may form two related eigenvector decompositions of $A$:
\[
  \highlight{A^TA = V \Lambda V^T, \qquad AA^T = U \Lambda U^T}.
\]
Since each eigenvector of $AA^T$ maps to an eigenvector of $A^TA$ and
vice-versa, we can be confident that the ranks of $AA^T$ and $A^TA$ are
identical. This is yet more justification of the assertion about the sizes of
$\Lambda_U$ and $\Lambda_V$ above.

\subsection{The singular value decomposition}

For the moment, suppose that there is some decomposition of an $N \times M$
matrix $A$ into a $N \times R$ orthogonal matrix $U$, a $M \times R$ matrix $V$
and some $R \times R$ matrix $\Sigma$ whose only non-zero terms are on the
diagonal:
\[
  \highlight{A = U \Sigma V^T}.
\]
Consider the form of the matrices $A^TA$ and $AA^T$:
\[
  A^TA = V \Sigma^T U U^T \Sigma V^T = V \Sigma^T \Sigma V^T, \qquad
  AA^T = U \Sigma V^T V \Sigma^T U = U \Sigma \Sigma^T U^T.
\]
If we make the observation that we may set $\Sigma^T \Sigma = \Lambda_V$ and
$\Sigma \Sigma^T = \Lambda_U$ then we have exactly the eigenvector
decomposition. Further, since $\Lambda_U = \Lambda_V = \Lambda$, then we can see
that \highlight{$\Sigma$ is the $R \times R$ diagonal matrix of eigenvalue
square-roots}.

The decomposition $A = U \Sigma V^T$ is called the \emph{singular value
decomposition} (SVD). \highlight{The columns of the matrices $U$ and $V$ are the
eigenvectors of $AA^T$ and $A^TA$ respectively} ordered by decreasing
eigenvalue. The non-zero elements of $\Sigma$ are called the singular values and
\highlight{the number of singular values is equal to the rank of $A$}.

\subsubsection{The geometric interpretation of the SVD}

We can interpret the SVD of a $N \times M$ matrix geometrically. The action of
the first orthonormal matrix $V^T$ is to transform a $M$-component vector into a
$R$-dimensional subspace. This will lose information if $R < M$. In this
subspace, the action of the matrix $A$ is to scale along each of the basis
vectors by the values of the diagonal of $\Sigma$. Finally, the orthonormal
matrix $U$ transforms from the $R$-dimensional subspace into the $U$-component
space we expect the result to be in. This transform cannot invent new
information; the result of applying $U$ is still $R$-dimensional, it is just a
$R$-dimensional subspace embedded in a $N$-dimensional space.

It is this `necking' into a $R$-dimensional space that shows up if the matrix
$A$ is invertible in and of itself. If $R < M$ information will be lost and so
the matrix could never be invertible.

\subsubsection{Matrix approximation}

\begin{figure}[p]
  \begin{centering}
    \begin{tabular}{cc}
      \includegraphics[width=0.4\columnwidth]{svd_original} &
      \includegraphics[width=0.4\columnwidth]{svd_50} \\
      Original (100\% storage) & 50\% of s.v.s (82\% storage)\\[1em]
      \includegraphics[width=0.4\columnwidth]{svd_25} &
      \includegraphics[width=0.4\columnwidth]{svd_10} \\
      25\% of s.v.s (41\% storage) & 10\% of s.v.s (16\% storage)\\[1em]
      \includegraphics[width=0.4\columnwidth]{svd_5} &
      \includegraphics[width=0.4\columnwidth]{svd_1} \\
      5\% of s.v.s (8\% storage) & 1\% of s.v.s (2\% storage)\\[1em]
    \end{tabular} \\
  \end{centering}
  \caption{\label{fig:svd-truncation}An example of truncating the singular
  values in the SVD. Under each image is shown the number of singular values
  (s.v.s) retained and the corresponding amount of storage required for the
  image as a fraction of the original.}
\end{figure}

The SVD may be used to approximate a matrix. Instead of retaining the full $R$
diagonal entries of $\Sigma$, we may truncate the matrix and keep only the first
$R' < R$ entries of $\Sigma$ and the first $R'$ columns of $U$ and $V$. Figure
\ref{fig:svd-truncation} gives an example of this so-called \emph{rank
reduction} applied to a matrix representing an image. The SVD can be used as a
na\"ive form of image compression; we managed to get to around 16\% of the
original image's storage by setting $R' = 0.1 R$.

\end{document}

% vim:sw=2:sts=2:autoindent:tw=80
